\begin{proposition}[E.24 --- Proved: the integral elliptic-involution constraint]
\label{E-integral-symmetry}
Retain the integral class $[\eta]\in H^1(G_{\mathbb Q,\{3,\infty\}},
A(-1))$ of E.19. Let $\iota$ be the automorphism
$[x:y:z]\mapsto[y:x:z]$ of the Fermat cubic with origin
$O=[1:-1:0]$. It induces an involution on the actual
based pro-three group, commuting with its arithmetic action.
Put $A^+=\ker(\iota-1:A\to A)$. The class $[\eta]$
is in the image of the specified inclusion
\[
 H^1(G_{\mathbb Q,\{3,\infty\}},A^+(-1))
 \longrightarrow H^1(G_{\mathbb Q,\{3,\infty\}},A(-1)),
 \tag{E24.1}
\]
and the source is
\[
 \mathbb Z_3^{81}\oplus
 (\mathbb Z/3)^{31}\oplus
 (\mathbb Z/9)^{12}\oplus
 (\mathbb Z/27)^4.\tag{E24.2}
\]
The inclusion in (E24.1) is split by the integral projector
$(1+\iota)/2$. This assertion constrains the genuine
integral class, including its torsion coordinates.
\end{proposition}
\begin{proof}
On the complete elliptic curve, $\iota=[-1]$.
At $O$ the parameter $z/x$ has derivative $-1$
under $\iota$. The compatible $3^r$th roots of $-1$
can all be chosen to equal $-1$. They are rational
and multiply to $1$ when the tangent change is
performed twice. The corresponding tangential
path therefore supplies an order-two based
automorphism commuting with $G_{\mathbb Q}$.
This constructs the claimed action, including the
comparison between its two tangent base points.

Use the ordered degree-one basis and the matrix $U$
of E.18. On cusp homology $\iota$ is the permutation
matrix $P$ with column permutation
$(3,4,5,0,1,2,7,6)$, in zero-based indexing. Its
matrix on $V$ is
\[
 I_V=\begin{pmatrix}
 -I_2&0\\(-U-PU)/3&P
 \end{pmatrix}.\tag{E24.3}
\]
Every entry is integral. The map is $-1$ on the
compact lift vectors $3f_i+U_i$, so it fixes
$w=(3f_1+U_1)\wedge(3f_2+U_2)$.
Let $I_A=\operatorname{Lie}_3(I_V)$.
It commutes with the Frobenius matrix $T$ and
the conjugation matrix $J$ of E.18.

Begin with any lift $\omega_0$ of $w$ to
$L=\gamma_2N/\gamma_4N$. Apply the commuting
projectors $(1+\iota)/2$ and $(1-c)/2$ to obtain
$\omega$ with $\iota\omega=\omega$ and
$c\omega=-\omega$. Its image in $M$ is still $w$.
For every arithmetic $g$ the formula
\[
 \eta(g)=\chi(g)^{-1}(g\omega-\chi(g)\omega)
\]
now has values in $A^+(-1)$, since $g$ commutes
with $\iota$. Replacing the lift changes $\eta$
by the actual coboundary of the difference of the
two lifts in $A(-1)$. Hence (E24.1) follows for
the original cohomology class.

The degree-three free Lie trace identity gives
$\operatorname{tr}(I_A)=-2$ and
$\operatorname{rank}A^+=164$. Also
$\operatorname{tr}(J|A^+)=-2$. In the free
pro-three presentation used in E.18--E.19, the
two generators are interchanged by $c$.
Consequently the same restriction-and-coboundary
calculation as in E.19 presents the group in
(E24.1) as the quotient of $A^+$ by the image of
$(T-7I)(1-J)$ on $A^+$; the factor $7$ is a unit
and is the cyclotomic twist in that calculation.

For an exact computation in the ambient integral
lattice use the matrix
\[
             (T-7I)(1-J)(1+I_A).
\]
Multiplication by $1+I_A$ has image $A^+$, since
two is a $3$-adic unit. The executed Hall-basis
program \texttt{class\_three\_symmetry.py} gives
the following elementary-divisor valuations:
\[
\begin{array}{c|rrrr}
\text{valuation}&0&1&2&3\\\hline
\text{multiplicity}&36&31&12&4.
\end{array}
\]
All 83 nonzero pivots are certified modulo $3^{10}$.
The exact rank cannot exceed
$(164-(-2))/2=83$, the dimension of the
$(-1)$-eigenspace of $J$ in $A^+$.
Thus these are all its nonzero elementary divisors.
The complement $A^-$ is a direct summand of the
ambient lattice and the relation image lies in
$A^+$; removing that free complement leaves
free rank $164-83=81$ and exactly the torsion
factors in (E24.2).

The rational averaging over the larger group in
E.23 has a denominator divisible by three.
Its comparison map is the rationalization of
the integral extension followed by the quotient
$A\to A/[V,w]$. The integral projector used here
has denominator two. The two conclusions are
therefore recorded along their respective maps;
neither calculation evaluates the remaining
arithmetic cocycle vector. The manuscript anchor
is \S51.16, folios 733v--736v, scans 172--178.
\end{proof}
